% =============================================================================
%  2AMD20 Knowledge Engineering --- Take-home Final Exam
%  Q4 2025-2026 :: TU/e
%
%  Margins follow \usepackage{fullpage} per the grading rubric.
%  Each answer lives in its own file under questions/; reorder via \input.
% =============================================================================

\documentclass[11pt,a4paper]{article}

% ---- Layout -----------------------------------------------------------------
% The grading rubric asks for "margins comparable to \usepackage{fullpage}".
% We use geometry with 1-inch margins (fullpage's effective layout) plus
% explicit headheight/headsep so fancyhdr's header does not collide with the
% first section title on each page.
\usepackage[a4paper, margin=1in, headheight=14pt, headsep=18pt, includeheadfoot]{geometry}
\usepackage[utf8]{inputenc}
\usepackage[T1]{fontenc}
\usepackage{lmodern}
\usepackage{microtype}
\usepackage[most]{tcolorbox}
\usepackage{tikz}
\usetikzlibrary{positioning, shapes.multipart, arrows.meta}
\newtcolorbox{topicquote}{
    colback=white,       % Light background tint
    colframe=white,       % No border line
    arc=0mm,              % Square corners
    halign=center,        % Center text inside the box
    valign=center,
    width=\textwidth,     % Take up full row width
    top=1em, bottom=1em   % Padding
}

% ---- Math / symbols ---------------------------------------------------------
\usepackage{amsmath, amssymb, amsthm}
\usepackage{mathtools}

% ---- Code / quoted blocks ---------------------------------------------------
\usepackage{listings}
\usepackage{xcolor}
\lstset{
    basicstyle=\ttfamily\small,
    breaklines=true,
    columns=fullflexible,
    frame=single,
    framerule=0pt,
    backgroundcolor=\color{black!4},
    keepspaces=true,
    showstringspaces=false,
}

% RDF / Turtle pseudo-syntax highlighting (optional)
\lstdefinelanguage{turtle}{
    morecomment=[l]{\#},
    morestring=[b]",
    sensitive=true,
    morekeywords={a, PREFIX, BASE},
}

% Cypher / GQL pseudo-syntax highlighting (optional)
\lstdefinelanguage{cypher}{
    morecomment=[l]{//},
    morestring=[b]",
    sensitive=false,
    morekeywords={MATCH, RETURN, WHERE, CREATE, MERGE, DELETE, SET, WITH,
                  OPTIONAL, UNWIND, AS, AND, OR, NOT, IN},
}

% ---- Graphics / hyperlinks --------------------------------------------------
\usepackage{graphicx}
\usepackage[colorlinks=true, linkcolor=black, urlcolor=blue, citecolor=blue]{hyperref}

% ---- Bibliography ------------------------------------------------------------
%  natbib + BibTeX is used because both ship with Tectonic and need no external
%  binary. If you prefer biblatex+biber, install biber first
%  (`brew install biber`), then swap these two lines for:
%     \usepackage[backend=biber,style=numeric-comp]{biblatex}
%     \addbibresource{references.bib}
%  and replace \bibliographystyle{...}\bibliography{references} below with
%  \printbibliography.
\usepackage[numbers,sort&compress]{natbib}

% ---- Headers ----------------------------------------------------------------
\usepackage{fancyhdr}
\pagestyle{fancy}
\fancyhf{}
\fancyhead[L]{2AMD20 Knowledge Engineering --- Take-home Final Exam}
\fancyhead[R]{\thepage}
\renewcommand{\headrulewidth}{0.4pt}

% ---- Custom question environment --------------------------------------------
% Usage:
%   \begin{question}{N}{Lecture number}{Short topic}
%     \questionbrief{prompt text}{referenced papers}{grading and length}
%     ... your answer ...
%   \end{question}
%
% Renders a labelled heading and a quoted brief block so the question text
% always travels with the answer in the rendered PDF.
\newenvironment{question}[3]
  {\clearpage\section*{Question #1 --- Lecture #2: #3}\addcontentsline{toc}{section}{Question #1 --- Lecture #2: #3}}
  {\vfill}

\newcommand{\questionbrief}[3]{%
  \begin{quote}\small
    \noindent\textbf{Prompt.}\quad #1\par\medskip
    \noindent\textbf{Referenced material.}\quad #2\par\medskip
    \noindent\textbf{Grading \& length.}\quad #3
  \end{quote}
  \par\noindent\rule{\linewidth}{0.4pt}\par\medskip
}

% =============================================================================
%  Front matter
% =============================================================================
\title{Take-home Final Exam \\
       \large 2AMD20 Knowledge Engineering, Q4 2025--2026}
\author{Adam Szelestey \\ Student ID: \texttt{2362783} \\ \texttt{a.szelestey@student.tue.nl}}
\date{\today}

% --- Custom Styling for Knowledge Engineering Exam ---
\usepackage{enumitem} % Allows for compact algorithms and lists

% Reusable, compact inline definition macro
% Usage: \kgedef{Number}{Title} Definition text starts here...
\newcommand{\kgedef}[2]{%
  \par\vspace{0.5em}\noindent\textbf{Definition #1 (#2):}\quad%
}

\begin{document}
\maketitle

\begin{abstract}
\noindent
This document contains my answers to five questions of the 2AMD20 take-home final exam,
each from a different lecture week, per the rules in the course syllabus.
\end{abstract}

\tableofcontents
\bigskip

% =============================================================================
%  Five answers, each from a different week. Each question lives in its own
%  file under questions/ to keep diffs small and review focused.
%  Reorder freely by reordering the \input lines.
% =============================================================================

% =============================================================================
%  Question 1 --- Lecture 2: Reification in property graphs
%  Source: docs/course_extracts/lecture2/06-exam-and-readings.md
% =============================================================================

\begin{question}{1}{2}{Reification in property graphs}

\questionbrief
  {%
    In this lecture we have discussed the notion of reification for RDF.
    Formulate a notion of reification in the world of property graphs and
    illustrate how reification in property graphs would work in practice using
    two small examples. One of those examples should be
    \emph{``The detective supposes that the butler killed the gardener''}.
    Since we are dealing with property graphs, make sure to include
    vertex/edge labels and key-value pairs in your proposal.%
  }
  {%
    No specific paper is required by the prompt. Useful background:
    Hogan et al., \emph{Knowledge graphs} (chs.~1--2);
    Angles \& Guti\'errez, \emph{Survey of graph database models} (ACM CSUR, 2008);
    Guarino, Sales \& Guizzardi, \emph{Reification and Truthmaking Patterns}
    (ER 2018).%
  }
  {%
    70\% completeness, correctness, and elegance of the proposal;
    30\% writing quality/clarity/completeness.
    Minimum length: \textbf{2 A4 full pages}.%
  }

% --- Your answer below -------------------------------------------------------

% TODO: define your notion of property-graph reification (motivation +
% formal/semi-formal definition), then provide the two required examples
% with explicit vertex/edge labels and key-value pairs.

\end{question}

% =============================================================================
%  Question 2 --- Lecture 3: Converting between RDF and MPG reification
%  Source: docs/course_extracts/lecture3/04-conclusion-and-readings.md
% =============================================================================

\begin{question}{2}{3}{Converting between RDF and MPG reification}

\questionbrief
  {%
    Sadoughi et al.\ have recently introduced \emph{Meta-Property Graphs}
    (MPG), a practical extension of Property Graphs. Give procedures for
    converting between reification in RDF data and reification in MPG data
    (\textbf{both directions}). Present your procedures in a self-contained,
    stand-alone write-up of your investigation.%
  }
  {%
    Primary: Sadoughi, Yakovets \& Fletcher, \emph{Breaking down the
    data-metadata barrier for effective Property Graph data management},
    EDBT 2025. Useful background: Guarino, Sales \& Guizzardi,
    \emph{Reification and Truthmaking Patterns} (ER 2018);
    Hogan et al., \emph{Knowledge graphs} (\S3.2).%
  }
  {%
    70\% quality and motivation/analysis of the proposal;
    30\% writing quality/clarity/completeness.
    Minimum length: \textbf{2 A4 full pages}
    (margins comparable to \texttt{\textbackslash usepackage\{fullpage\}}).%
  }

\section{Converting between RDF and MPG reification}
% --- Your answer below -------------------------------------------------------

% TODO:
%   (a) motivate the conversion problem;
%   (b) RDF reification -> MPG procedure (with worked example);
%   (c) MPG reification -> RDF procedure (with worked example);
%   (d) brief discussion of correctness / round-tripping limits.
%
% NOTE: keep this answer substantively different from Question 3 (L8).
% L8 asks for a one-direction direct mapping; here you must cover BOTH
% directions and justify the design choices, not just exhibit a mapping.
\subsection{Theoretical Foundation \& Problem Statement}

Knowledge Graphs 

\subsection{RDF to MPG reification}

\subsection{MPG to RDF reification}

\end{question}

% =============================================================================
%  Question 3 --- Lecture 8: Direct mapping from reified RDF to reified MPG
%  Source: docs/course_extracts/lecture8/06-readings-and-exam.md
% =============================================================================

\begin{question}{3}{8}{Direct mapping from reified RDF to reified MPG}

\section{Direct mapping from reified RDF to reified MPG}
\subsection{Introduction and Theoretical Foundation}

To establish a direct mapping from standard RDF reification to Meta-Property Graphs (MPG), we must first formally define the structural characteristics of both the input space (RDF) and the target output space (MPG).

\kgedef{1}{Standard RDF Reification (Input)} Let $G_{RDF}$ be an RDF graph containing a reified statement. In standard RDF reification, a statement is structurally represented by an identifier node $i$ of type \texttt{rdf:Statement}. This node is linked to the base elements via three core triples: $(i, \texttt{rdf:subject}, s)$, $(i, \texttt{rdf:predicate}, p)$, and $(i, \texttt{rdf:object}, o)$. Metadata attached to this statement is represented by additional triples. To ensure completeness during mapping, we distinguish between two types of metadata triples: \textit{data properties} $(i, p_d, v_d)$ where the target is a literal value, and \textit{object properties} $(i, p_o, n_o)$ where the target is another entity \cite{hernandez_reifying_nodate, noauthor_rdf_nodate, hogan_knowledge_2022}.

\kgedef{2}{Meta-Property Graph (Output)} The target database is a Meta-Property Graph formally defined as the tuple $G_{MPG} = (N, E, P, L, \lambda, \mu, \sigma, \upsilon, \eta, \rho)$. In this structure, $N$ is the set of nodes, $E$ is the set of edges (including directed edges $E_d$), $P$ is the set of properties, and $L$ is the set of label identifiers. The model uses specific mapping functions: $\eta$ for edge topology, $\lambda$ and $\mu$ for labeling, $\sigma$ and $\upsilon$ for property assignment. Crucially, it introduces the $\rho$ function, which natively assigns a set of objects to a specific node to encapsulate a sub-structure \cite{sadoughi_meta-property_2025}.

\subsection{Direct Mapping Algorithm: RDF Reification to MPG Reification}

The following pseudo-code details the transformation logic to map an RDF reified statement into an MPG reified sub-structure. 

\vspace{0.5em}
\noindent\hrulefill
\begin{flushleft}
\textbf{Algorithm:} Map\_RDF\_to\_MPG($G_{RDF}$) \\
\textbf{Input:} An RDF Graph $G_{RDF}$ \\
\textbf{Output:} A Meta-Property Graph $G_{MPG} = (N, E, P, L, \lambda, \mu, \sigma, \upsilon, \eta, \rho)$
\end{flushleft}
\vspace{-1.2em}
\noindent\hrulefill
\begin{enumerate}[leftmargin=*, noitemsep, topsep=4pt]
    \item \textbf{For} every statement node $i \in G_{RDF}$ of type \texttt{rdf:Statement} \textbf{do}:
    \item \quad Extract base triples: $(i, \texttt{rdf:subject}, s)$, $(i, \texttt{rdf:predicate}, p)$, and $(i, \texttt{rdf:object}, o)$
    \item \quad \textit{// Step 1: Map Base Nodes}
    \item \quad Add $n_s, n_o$ to $N$
    \item \quad \textit{// Step 2: Map Relationship Predicate}
    \item \quad Create directed edge $e \in E_d$
    \item \quad Set $\eta_s(e) \leftarrow n_s$ and $\eta_t(e) \leftarrow n_o$
    \item \quad Set $p \in \mu(\lambda(e))$
    \item \quad \textit{// Step 3: Create Meta-Node and Apply $\rho$}
    \item \quad Create meta-node $n_i \in N$
    \item \quad Assign $\rho(n_i) \leftarrow \{n_s, e, n_o\}$
    \item \quad \textit{// Step 4: Attach Metadata}
    \item \quad \textbf{For} every metadata triple $(i, p_m, target) \in G_{RDF}$ \textbf{do}:
    \item \quad \quad \textbf{If} $target$ is a literal value $v_d$ \textbf{then}
    \item \quad \quad \quad Create property $k \in P$ with $\upsilon(k) \leftarrow (p_m, v_d)$
    \item \quad \quad \quad Add $k$ to $\sigma(n_i)$
    \item \quad \quad \textbf{Else if} $target$ is an entity $n_o'$ \textbf{then}
    \item \quad \quad \quad Create directed edge $e_o \in E_d$
    \item \quad \quad \quad Set $\eta_s(e_o) \leftarrow n_i$ and $\eta_t(e_o) \leftarrow n_o'$
    \item \quad \quad \quad Set $p_m \in \mu(\lambda(e_o))$
    \item \quad \quad \textbf{End If}
    \item \quad \textbf{End For}
    \item \textbf{End For}
\end{enumerate}
\noindent\hrulefill
\vspace{0.5em}

\subsection{Explanation of the Mapping Steps}

\begin{itemize}[leftmargin=*]
    \item \textbf{Step 1 (Map the Base Nodes):} We begin by identifying the statement node $i$ in $G_{RDF}$. We extract the targets of its \texttt{rdf:subject} and \texttt{rdf:object} triples, representing $s$ and $o$. We then map $s$ and $o$ to distinct nodes $n_s$ and $n_o$, and add them to the MPG node set, such that $n_s, n_o \in N$ \cite{hernandez_reifying_nodate, sadoughi_meta-property_2025}.
    
    \item \textbf{Step 2 (Map the Relationship Predicate):} Next, we extract the target of the \texttt{rdf:predicate} triple, representing $p$. We create a new directed edge $e$ in the MPG edge set, such that $e \in E_d$. We use the topology function $\eta$ to assign its endpoints: $\eta_s(e) = n_s$ and $\eta_t(e) = n_o$. Finally, we use the labeling functions $\lambda$ and $\mu$ to map the predicate name to the edge, such that $p \in \mu(\lambda(e))$ \cite{sadoughi_meta-property_2025}.
    
    \item \textbf{Step 3 (Create the Meta-Node and Apply $\rho$):} To cleanly map the RDF statement identifier $i$, we create a new meta-node $n_i$ in the MPG node set ($n_i \in N$). Instead of creating structural helper edges (which leads to graph bloat in standard RDF), we use the native $\rho$ (rho) function to encapsulate the base edge into a reified sub-structure by assigning $\rho(n_i) = \{n_s, e, n_o\}$. This formally makes $n_i$ a first-class citizen representing the original relationship \cite{sadoughi_meta-property_2025}.
    
    \item \textbf{Step 4 (Attach the Metadata):} Finally, we scan $G_{RDF}$ for any remaining triples where the statement identifier $i$ is the subject, mapping them to $n_i$ depending on their target type:
    \begin{itemize}
        \item \textit{For Data Properties} (metadata with literal values $v_d$): We create a new property object $k$ in the property set ($k \in P$). We use the $\upsilon$ (upsilon) function to assign the key-value pair: $\upsilon(k) = (p_d, v_d)$. We then attach this property directly to the meta-node using the $\sigma$ (sigma) function by adding it to the node's compatible property set: $k \in \sigma(n_i)$ \cite{sadoughi_meta-property_2025}.
        \item \textit{For Object Properties} (metadata pointing to other entities $n_o'$): Because $n_i$ is a recognized structural node in the MPG, we create a new directed edge $e_o \in E_d$. We use the $\eta$ function to define its topology originating from the meta-node: $\eta_s(e_o) = n_i$ and $\eta_t(e_o) = n_o'$. We map the predicate $p_o$ to the edge label such that $p_o \in \mu(\lambda(e_o))$ \cite{sadoughi_meta-property_2025, hogan_knowledge_2022}.
    \end{itemize}
\end{itemize}

\end{question}
\begin{question}{4}{5}{Solving knowledge gap bridge with relational database}

% \questionbrief
%   {%
%     Consider again the E-commerce use case of Sequeda et al. discussed in 
% this lecture. In this use case, the gap between the data producers and data 
% consumers was bridged using a knowledge graph. Is it possible to bridge 
% this gap using a relational database instead of a knowledge graph? Provide 
% arguments that support your answer.%
%   }
%   {%
%     Sequede et al. \cite{sequeda_pay-as-you-go_2019}%
%   }
%   {%
%     Grading: 70\% for quality and and depth of your analysis, 30\% for quality/
%     clarity/completeness of writing.  Your answer to this question should be at 
%     least one A4 fullpages, i.e., page margins equivalent to%
%   }
\section{Solving knowledge gap bridge with relational database}

\subsection{Introduction}

Yes, you can close the knowledge gap in other names data-meaning gap in the discusses E-Commerce use case with Relational Database and no Knowledge Graphs. In more detail the procedure can be done utilizing relational database mappings in other names Ontology-Based Data Access (ODBA) \cite{poggi_linking_2008} methods. This is done by following the same process proposed by \citet{sequeda_pay-as-you-go_2019} but replacing the underlying knowledge graph architecture with an equivalent ODBA system.
To fully understand the question, it is important to have clear definition of data-meaning gap: based on \citet{sequeda_pay-as-you-go_2019} data-meaning gap occurs when data consumers cannot answer there own questions without heavy IT support. It is also essential to have a shared understanding of relational databases, so here it will be defined as database architectures organizing their data in tables with strict schemes \cite{codd_relational_1970}.

\subsection{Premises}

As \citet{sequeda_pay-as-you-go_2019} proved in their work the data-meaning gap can be closed with the usage of the \textit{Pay-as-you-go} Knowledge Graph creation procedure. They proved it by implementing it in real commercial cases where they observed satisfactory results.

It is important to note that the real gap is human not technical. The step that explicitly depends on knowledge graphs is the . Where you can replace the underlying knowledge graph with in its expressiveness equal relational database. This is best done via ODBA methods \cite{poggi_linking_2008} with which we can keep the rich expressiveness of KGs, but effectively rely on relational databases. Ontology-Based Data Access (OBDA) is a method that allows users to query pre-existing, complex databases using a simplified, high-level conceptual vocabulary (an ontology) rather than the database's underlying technical schema.

In short Knowledge Graph is just the architecture that might happen to be more efficient in few cases and make future work more consistent and flexible, but the same can be achieved through relational databases \cite{sakr_relational_2010}) we would get by solving phase 3. Actually, many knowledge graphs are implemented with relational databases.
A valid concern here is that the main advantage of knowledge graphs is their flexibility, which in fact is lost partly with ontology based methods.

\subsection{Conclusion}

In conclusion, as pay-as-you go solves the data-meaning gap and we can create an equivalent method with relational databases it is proven that the data-meaning gap could be solved with only relational databases and no knowledge graphs. Consequently, the architecture of the database in most knowledge gap problems is not the most concerning problem, but rather defining the mediator ontology that satisfies most requirements.

% --- Your answer below -------------------------------------------------------
\end{question}

% TODO: pick two more questions from a different week each.
%   Candidates (see docs/course_extracts/exam.md):
%   - Lecture 4: Containment of CQ_k queries
%   - Lecture 5: Relational DB vs. KG for the e-commerce use case (>= 1 page)
%   - Lecture 6: Review of a "data context" paper (>= 2 pages)
%   - Lecture 7: Cameron's Laws of Identity and societal challenges (1-2 pages)
%
% Once chosen, create the file under questions/ following the existing pattern
% and add \input{questions/q4_...} and \input{questions/q5_...} below.
%
% \input{questions/q4_...}
% \input{questions/q5_...}

% =============================================================================
%  AI usage acknowledgement (required by the course AI policy)
% =============================================================================
\clearpage
\section*{AI usage acknowledgement}
\addcontentsline{toc}{section}{AI usage acknowledgement}

Per the course AI policy, all writing in this document was produced by me.
AI tools were used for the following purposes:

\begin{itemize}
    \item \textbf{Tool:} NotebookLM \begin{itemize}
        \item \textbf{Purpose:} Help in understanding collected research material.
        \item \textbf{Prompts / scope:} Prompts here were extensive and covered multiple topics. Full excerpt can be shared.
    \end{itemize}
    \item \textbf{Tool:} Gemini \begin{itemize}
        \item \textbf{Purpose:} Formatting latex, revising text, searching for research material
        \item \textbf{Prompts / scope:} 
    \end{itemize}
\end{itemize}

% =============================================================================
%  References
%
%  Cite work in your answers with \citep{key} / \citet{key}, where `key` is
%  any entry from references.bib. While no \cite command is present the
%  References section will render empty (this is expected; BibTeX will warn
%  "I found no \citation commands" but the build still succeeds).
%
%  Tip: while drafting, add \nocite{*} just above \bibliography to dump every
%  entry from references.bib into the References section so you can see them.
% =============================================================================
\bibliographystyle{plainnat}
\bibliography{references}

\end{document}
